\section{Introduction}
\IEEEPARstart{F}{unction} computation over multiple-access channels uses the superposition of simultaneous transmissions to evaluate a desired function of distributed observations at a common receiver~\cite{nazer2007,sahin2023survey}. In the conventional analog formulation of over-the-air computation, the observations and the desired function are real-valued. Each transmitter preprocesses its observation, and the receiver applies a postprocessing function to the sum; analog amplitude modulation provides a direct implementation of this principle~\cite{goldenbaum2015,channelcomp2024}. For example, transmitting scaled observations allows the receiver to compute their arithmetic mean. The complex baseband representation of the wireless channel specifies two real signal coordinates and does not require a complex-valued desired function.

Our starting point is function-dependent digital modulation, as developed in ChannelComp~\cite{channelcomp2024}. A digital encoder assigns a constellation point to each quantized observation. Because the channel adds the transmitted symbols, their placement determines which input configurations the receiver can distinguish. For computation, the relevant distances are between received sums associated with different desired function values. Two such sums can be close even when the corresponding function values differ substantially, allowing a small disturbance to cause a large computation error. We therefore seek an input-to-symbol mapping that enables the desired computation and controls its sensitivity to noise under limited transmit power.

We study this design question for real-valued observations through a shared continuous map $\enc:I\to\C$, where $I=[-1,1]\subset\R$. As the observation varies, its transmitted symbol traces a curve in the in-phase and quadrature plane; restricting the map to a finite input alphabet gives a digital constellation. The two coordinates provide two additive signal components within one complex channel use. For variance, a symbol proportional to $s+i s^2$ produces a received sum proportional to $\sum_k s_k+i\sum_k s_k^2$. The receiver can therefore recover the first moment $m$ and second raw moment $q$ and compute $q-m^2$. This example explains the role of the complex-plane curve: it specifies how to convey the quantities needed for a real-valued computation through the two available signal coordinates.

A direct transmission of $s$ and a prescribed preprocessing function $\phi(s)$ establishes achievability, but leaves the optimal encoder geometry undetermined. Increasing the scale of one transmitted quantity can reduce its sensitivity to noise, while the pointwise power constraint limits the scale available to both quantities. We must determine whether exact computation forces every continuous encoder to preserve these two sums, and then optimize their scaling, relative orientation, and common offset. The continuous formulation lets us address both questions over all shared continuous encoders on $I$. The input parameterization remains part of the design: two encoders with the same curve image can produce different received sums for the same observations.

We address these questions for real-valued desired functions $f_{\phi,\psi}:I^K\to\R$ of the form
$$
 f_{\phi,\psi}(\mathbf s)
 =\frac1K\sum_{k=1}^{K}\phi(s_k)
 -\psi\!\left(\frac1K\sum_{k=1}^{K}s_k\right),
 \qquad \mathbf s\in I^K,
$$
where $\phi,\psi:I\to\R$ are prescribed. This class includes variance when $\phi(s)=\psi(s)=s^2$. Section~\ref{sec:model} gives further examples and explains its relation to scalar nomographic representations. We use one shared continuous encoder and one complex channel use. Our design objective is the infimum, over input pairs with different desired function values, of the squared distance between their received sums divided by their squared function difference. The inverse square root of this \emph{squared-separation score} is the smallest uniform disturbance-amplification constant achievable by an exact decoder.

\subsection{Contributions}
We connect exact computation, optimal encoder geometry, and receiver error through the following results.
\begin{enumerate}
\item \emph{Necessary encoder structure.} For $K\ge4$ under the rigidity assumptions in Table~\ref{tab:assumptions}, we prove that every continuous encoder computing $f_{\phi,\psi}$ exactly has the form $\enc(s)=c+us+v\phi(s)$, where $c,u,v\in\C$ and $u,v$ are linearly independent over $\R$. Thus every exact encoder preserves the two sums $S=\sum_k s_k$ and $T=\sum_k\phi(s_k)$. We derive this conclusion from input triples with equal received sums without assuming that the encoder is differentiable. The result permits preprocessing functions $\phi$ with affine subintervals.
\item \emph{Optimal encoding and sharp disturbance bounds.} Under the stronger assumptions in Table~\ref{tab:assumptions}, we express the exact squared-separation score through two endpoint gradients and obtain convex formulations of the peak-constrained encoder design. We derive an equivalent enclosing-ellipse representation of the curve $(s,\phi(s))\in\R^2$. A unified theorem gives explicit optimal encoders for quadratic and quartic preprocessing. We extend the encoder characterization and convex design to functions of two prescribed additive statistics under a common-completion identifiability condition. Polynomial moment interactions, including $mq$, and a nonpolynomial transformation of variance illustrate the extension. For an exponential preprocessing function, we certify a numerical design with a relative global gap below $0.006\%$. A decoder using a fixed linear transformation, scalar clipping, and function evaluation attains the smallest uniform disturbance-amplification constant, including when $u$ and $v$ are nonorthogonal.
\item \emph{A global robustness--accuracy tradeoff.} For a specified input distribution and the constructed decoder, we minimize the leading mean-squared error (MSE) coefficient under circular complex Gaussian noise over all continuous exact encoders meeting a required separation score. A five-variable convex program attains the optimum and determines a unique $2\times2$ Gram matrix of the encoder coefficients $u,v$. The optimal value is continuous and convex in the required score, with an initial plateau followed by a strictly increasing branch. For $\phi(s)=s^2$, we obtain an explicit scalar-root characterization. These results quantify how the optimal curve changes when the designer varies the required uniform robustness and optimizes average small-noise accuracy.
\end{enumerate}
Paired finite-noise experiments, deterministic implementation-error bounds, and a qualitative approximate-rigidity result support the analysis. We provide the detailed optimality proofs and comparisons with other pairwise design criteria in the appendices.

\subsection{Related Work}
Nazer and Gastpar established an information-theoretic framework for computation over multiple-access channels~\cite{nazer2007}. Their compute-and-forward construction uses nested lattice codes to recover linear combinations of transmitted messages~\cite{nazer2011}. For nonlinear real-valued functions, nomographic computation combines local preprocessing, channel summation, and scalar postprocessing. Goldenbaum, Boche, and Sta\'nczak combine this representation with nested lattice coding and analyze computation rates under a block power constraint~\cite{goldenbaum2015}. Our desired function depends on two prescribed real sums, which we encode jointly in the two coordinates of a single complex symbol. We characterize the necessary continuous encoder structure and optimize its sensitivity to disturbances under a pointwise peak constraint.

ChannelComp formulates function-dependent modulation through separation constraints on received superpositions~\cite{channelcomp2024}. Extensions and variants address multiple-antenna computation, blind federated learning, structured digital coding, and constellation sampling~\cite{veccomp2026,blindfl2024,sumcomp2025,pyramidsampling2025}. The original finite-input squared-distance formulation motivates our objective, but constrains the norm of a finite modulation vector rather than each symbol's energy. Noise-aware ChannelComp uses criteria adapted to the noise distribution~\cite{noiseaware2025}; its Gaussian criterion has an uncentered exponential numerator. Analytical optimization also exists for prescribed sum-computing grid constellations with ML or MAP decoding~\cite{towards2025}. Our continuous curve formulation complements these digital designs by establishing an encoder converse and a global squared-separation optimum for the stated classes of nonlinear real-valued functions.

Shared sum representations also underlie Deep Sets~\cite{zaheer2017}. Wagstaff \emph{et al.}\ establish latent-dimension requirements for universal continuous representation and approximation~\cite{wagstaff19,wagstaff22}. Their exact-max argument combines the analysis of equal encoded sums with topology; in particular, \cite[Lemma 17]{wagstaff22} treats a fixed cardinality. These results allow particular desired functions to admit smaller representations. Chen, Chen, and Lu establish representation through multisymmetric polynomial generators and examine singularities of the postprocessing function~\cite{chen2024}.

The closest affine-uniqueness precedent is the restricted Cauchy--Pexider theorem of Charalambides~\cite[Theorem 1.6]{charalambides2013}. For smooth curves without affine subarcs, it identifies affine summands from an additive factorization; its proof regularizes the summands before differentiation. We begin with exact computation of a scalar function and derive the required local factorization from the implication that equal received sums give equal function values. Our assumptions also permit $C^1$ curves with affine subintervals. We retain a self-contained proof of this reduction and do not claim affine uniqueness as a new general principle. Scalar monotone-generator uniqueness for quasi-arithmetic means provides a further classical precedent~\cite[Theorems A and D]{matkowski2015}.

The small-noise objective has the weighted inverse-trace form familiar from A-optimal experimental design. Conic formulations for such criteria and additional resource constraints are classical~\cite{sagnol2010}; semidefinite methods also cover continuous polynomial and rational design spaces~\cite{papp2011}. We apply these tools after identifying the encoder family forced by exact computation. Our feasible matrices arise from fitting the prescribed curve $(s,\phi(s))$ into a peak-power disk through an affine transformation. The resulting robustness--accuracy tradeoff is global for this computation class and the stated decoder and small-noise model. Matrix optimization and enclosing-ellipsoid principles themselves are established~\cite{boyd2004}.

We use classical invariance of domain~\cite{hatcher} and scalar Lipschitz extension~\cite{mcshane1934} to establish the encoder characterization and its operational interpretation.

Section~\ref{sec:model} introduces the channel, desired function class, and objective. Section~\ref{sec:design} combines the necessary encoder structure with optimal peak-constrained design. Section~\ref{sec:examples} presents the quadratic and quartic cases and the extension to two prescribed additive statistics; Section~\ref{sec:decoding} constructs the associated receivers. Section~\ref{sec:pareto} develops the small-noise tradeoff. Sections~\ref{sec:certified} and~\ref{sec:evidence} address certified numerical design and performance. The appendices contain the detailed proofs, additional consequences, and Lean-verification scope.
